% Chapter 1

\chapter{Introduction}
\label{ch:introduction}

Autonomous robotic systems are increasingly being deployed in environments where direct human operation is costly, dangerous, or operationally infeasible. \citeauthor{bhatia2026integration} show that unmanned and intelligent systems are becoming increasingly important in defense and other safety-critical domains \citep{bhatia2026integration}. Within this broader field, Unmanned Ground Vehicles (UGVs) play a particularly important role because they can navigate terrain that is inaccessible or hazardous to human operators while carrying mission-specific sensing, computation, and payload systems, as discussed by \citeauthor{krecht2023resilience} \citep{krecht2023resilience}. The wider relevance of such systems is also visible in ground-disaster-response research, where \citeauthor{sotovergel2023transforming} emphasize the growing need for rapid and reliable robotic operation \citep{sotovergel2023transforming}.

Despite this potential, robust UGV navigation remains a difficult problem. Ground vehicles operate with a limited viewpoint and must interpret terrain, obstacles, and path feasibility from sensors mounted close to the ground. This creates uncertainty in rough outdoor environments, where local measurements alone may be insufficient for reliable long-range maneuver decisions. \citeauthor{chen2021cnn} identify vision-proprioception terrain understanding as an important direction for improving UGV robustness \citep{chen2021cnn}. At the same time, \citeauthor{chen2026obstacle} show that complex outdoor environments still challenge reliable autonomous movement even when perception quality improves \citep{chen2026obstacle}. As a result, the challenge is not only to avoid immediate collisions, but also to infer how the vehicle should move through a partially observable environment in a safe and purposeful way.

One promising way to reduce this limitation is to combine the local ground perspective of a UGV with the broader aerial perspective of an Unmanned Aerial Vehicle (UAV). \citeauthor{asadi2020integrated} demonstrate the practical value of integrated UAV--UGV cooperation for sensing and data collection in challenging environments \citep{asadi2020integrated}. More recent work by \citeauthor{mondal2025risk} also shows growing interest in coordinated UAV--UGV decision-making under operational constraints \citep{mondal2025risk}. In such a cooperative setup, the UAV can act as a mobile scout that improves the information available to the UGV, enabling more informed maneuver decisions than would be possible from ground sensing alone.

At the same time, recent advances in learning-based navigation suggest that trajectory prediction can provide an alternative to purely handcrafted maneuver rules. \citeauthor{alahi2016sociallstm} showed that recurrent models such as LSTMs can infer future motion from sequential history \citep{alahi2016sociallstm}. \citeauthor{chen2024data} further support the relevance of sequence-based prediction for ground navigation through data-driven trajectory planning for UGVs \citep{chen2024data}. \citeauthor{vaswani2023attention} introduced the attention mechanism that later became central to Transformer-based temporal modeling \citep{vaswani2023attention}. \citeauthor{giuliari2020transformer} then transferred these ideas directly to trajectory forecasting \citep{giuliari2020transformer}. Likewise, \citeauthor{kipf2017semi} showed that graph-based models can represent relational dependencies between multiple entities \citep{kipf2017semi}. \citeauthor{yu2018stgcn} further demonstrated how such relations can be modeled jointly with temporal evolution \citep{yu2018stgcn}. These developments make it possible to model navigation not simply as reactive control, but as a structured prediction problem in which future movement is inferred from past motion, local terrain cues, and multi-agent context.

However, translating these ideas into a practical UAV--UGV navigation framework introduces several open challenges. First, the UGV must still remain sensitive to terrain and obstacle structure rather than relying on abstract trajectory prediction alone. Second, the interaction between UGV and UAV must be represented in a way that preserves useful relational information without making the model unnecessarily complex. Third, any cooperative navigation framework must account for the fact that information exchanged between agents is not always ideal. In realistic deployments, communication is affected by delay, jitter, and packet loss; therefore, the benefit of UAV-assisted guidance depends not only on perception quality, but also on the robustness of the communication channel. If such effects are ignored during development, the resulting performance in simulation may overestimate what can be expected in more realistic conditions.

This thesis addresses these challenges by developing and evaluating a cooperative UAV--UGV navigation framework that evolves from a rule-based maneuvering baseline into a learning-based trajectory prediction system. A rule-based controller is first designed to generate navigation behavior, safety transitions, and cooperative scouting logic in simulation. The resulting trajectories are then transformed into a structured training dataset, from which learning-based models are trained to predict short-horizon UGV trajectories. These models are subsequently evaluated both offline, using prediction metrics, and online, by reinserting trained model weights into live simulation. To extend the realism of the evaluation, an OMNeT++ communication layer is incorporated to model degraded information exchange between aerial and ground agents, following the communication-simulation direction established by \citeauthor{varga2008omnet} \citep{varga2008omnet}.

The significance of this work lies in the way it connects several normally separate threads: rule-based robotic navigation, trajectory learning, UAV-assisted perception, and communication-aware evaluation. Rather than treating learning-based prediction as a purely abstract machine learning task, the thesis places it within a complete operational framework where sensing, control, coordination, and communication all influence the final maneuvering behavior. This makes the work relevant not only from a modeling perspective, but also from the standpoint of system-level autonomous robotics.

\section{Background and Motivation}

The motivation for this thesis arises from the need for autonomous systems that can continue to operate effectively in uncertain and high-risk environments. In many real-world scenarios, a UGV must reason about more than a single obstacle directly in front of it; it must also consider terrain continuity, possible route quality, and whether external support from cooperating agents can improve decision quality. Traditional reactive methods remain useful because they are interpretable and fast, but they often become brittle when the environment is irregular, partially observable, or dynamically changing. In contrast, learning-based models offer the possibility of inferring more general maneuver patterns from data, especially when these data encode the relationships between past motion, surrounding context, and successful future actions.

The UAV--UGV setting is especially motivating because it naturally introduces complementary viewpoints and information asymmetry. The UGV provides direct interaction with the terrain but suffers from limited look-ahead perception, whereas the UAV can observe broader spatial structure but depends on communication and coordinated interpretation. This thesis is motivated by the question of whether learning-based trajectory prediction can make better use of that cooperative structure than purely rule-based systems, while still remaining practical for system-level simulation and future deployment-oriented evaluation.

\section{Problem Statement}

The specific problem addressed in this thesis is how to design a cooperative navigation framework in which a UGV can predict and execute effective short-horizon maneuvers toward a goal while benefiting from UAV-assisted context under realistic communication conditions. Existing rule-based methods can provide robust safety behavior and interpretable decision logic, but they may become indecisive or overly local in complex terrain. Existing learning-based methods can improve predictive capability, but they often assume simplified perception, single-agent settings, or ideal information exchange.

Accordingly, the problem is not only to predict a trajectory, but to do so in a way that remains meaningful in a multi-agent robotic setting where terrain, obstacle structure, and communication quality all matter. The central challenge is therefore to bridge system-level robotic navigation and data-driven trajectory learning within one coherent experimental framework.

\section{Research Aim}

The aim of this thesis is to develop and evaluate a cooperative UAV--UGV navigation framework in which learning-based trajectory prediction is trained from rule-based maneuver data and assessed under realistic simulation and communication conditions.

To achieve this aim, the thesis pursues the following objectives:

\begin{enumerate}
    \item Develop a rule-based UAV--UGV navigation baseline capable of goal-directed movement, obstacle handling, terrain-aware support, and cooperative scouting in simulation.
    \item Record and transform simulation runs into a structured dataset suitable for supervised trajectory prediction.
    \item Design and train learning-based models that use sequential motion history, compact scan-derived features, and multi-agent relational context to predict short-horizon UGV trajectories.
    \item Compare alternative model families to determine which architecture provides the best balance of predictive accuracy and operational usefulness.
    \item Reinsert trained model weights into live simulation and compare learned-controller behavior against the rule-based baseline.
    \item Evaluate the effect of communication degradation on cooperative navigation by integrating OMNeT++ into the simulation workflow.
    \item Measure runtime and resource usage for both rule-based and learning-based execution in order to support later discussion of deployment-oriented feasibility.
\end{enumerate}

\section{Research Questions}

\textbf{Primary Research Question}

\textbf{PRQ1:}
How effectively can a learning-based trajectory prediction framework improve cooperative UAV--UGV maneuver generation in comparison to a rule-based baseline under terrain uncertainty and communication-aware simulation conditions?

\bigskip

\textbf{Secondary Research Questions}

\textbf{SRQ1: Model Design} \\
How do sequential, graph-based, and attention-based trajectory prediction models differ in their ability to represent UGV motion, UAV context, and short-horizon maneuver intent?

\medskip

\textbf{SRQ2: Cooperative Context} \\
To what extent does incorporating UAV-related relational information improve predicted UGV trajectory quality compared with models that rely only on ego-motion and local scan-derived features?

\medskip

\textbf{SRQ3: System-Level Performance} \\
How does the best-performing learned model behave when deployed back into live simulation, and how does this behavior compare with the original rule-based maneuvering framework?

\medskip

\textbf{SRQ4: Communication Awareness} \\
How sensitive is the cooperative navigation framework to degraded communication conditions such as delay, jitter, and packet loss when UAV-assisted guidance is transmitted through an OMNeT++ communication layer?

\medskip

\textbf{SRQ5: Runtime Feasibility} \\
What are the observed runtime and resource requirements of the rule-based and learning-based navigation pipelines, and what do they suggest about future deployment-oriented use on constrained systems?

\section{Scope and Delimitations}

This thesis is limited to a simulation-based study carried out in a ROS~2 and Gazebo environment with an OMNeT++ communication layer used for network-aware evaluation. The work focuses on one UGV and cooperative UAV support within a controlled simulation world rather than on full physical deployment. The learning models predict short-horizon future trajectories in the UGV local frame rather than issuing direct low-level actuator commands. The thesis also focuses on compact scan-derived and relational features rather than end-to-end learning from raw camera imagery or dense point clouds. These delimitations are intentional and help keep the study focused on the interplay between navigation logic, trajectory prediction, multi-agent context, and communication robustness.

\section{Thesis Contributions}

The main contributions of this thesis are as follows:

\begin{enumerate}
    \item A rule-based cooperative UAV--UGV maneuvering framework that integrates local obstacle handling, terrain-oriented support, and recovery behavior in simulation.
    \item A data-generation pipeline that converts rule-based simulation runs into structured trajectory-learning samples with shared train, validation, and test splits.
    \item A comparative learning framework for short-horizon UGV trajectory prediction using sequential, graph-based, and Transformer-enhanced model variants.
    \item A live simulation evaluation pipeline that deploys trained trajectory models back into the robotic simulation loop for system-level comparison.
    \item A communication-aware evaluation setup that integrates OMNeT++ into the cooperative robotics workflow to study the effect of degraded information exchange.
    \item A runtime and resource measurement workflow that supports comparison between rule-based and learning-based navigation from a deployment-oriented perspective.
\end{enumerate}

\section{Thesis Structure}

The remainder of this thesis is organized as follows. Chapter~2 presents the literature review and discusses prior work on UGV navigation in unstructured terrain, UAV-assisted navigation, rule-based maneuvering, learning-based trajectory prediction, graph-based multi-agent reasoning, and communication-aware robotic coordination. Chapter~3 describes the system architecture and experimental framework, including the robotic platforms, sensor setup, ROS~2 and Gazebo integration, mission scenario, and system-level information flow. Chapter~4 explains the methodology, beginning with the rule-based maneuvering baseline and continuing through dataset generation, model design, training procedure, evaluation strategy, performance metrics, runtime measurement, and communication-aware assessment. Chapter~5 reports the experimental results, covering rule-based behavior, dataset generation outcomes, offline model scores, live simulation performance, and runtime/resource comparisons. Chapter~6 discusses the findings in relation to the research questions, including the role of model components, the connection between offline and online behavior, edge-device implications, communication effects, and study limitations. Finally, Chapter~7 concludes the thesis by summarizing the main findings, revisiting the research questions, and outlining practical implications and recommendations for future work.
